\chapter{Trigonometria Contextualizada}

\begin{planningbox}{Bloque C05}
Los problemas de este capitulo obligan a controlar unidades, cuadrantes, lectura geometrica y
eleccion de herramienta. La dificultad no esta solo en calcular, sino en decidir si conviene
convertir, descomponer un triangulo, usar una identidad o resolver una ecuacion.
\end{planningbox}

\section{[C05.S01] Grados y radianes}

\begin{exercise}{[CTX-C05.S01-01] Giro acumulado de una camara panoramica}
Una camara de vigilancia parte orientada a \(35^\circ\). Durante una prueba automatica ejecuta
tres ordenes consecutivas:
\begin{enumerate}
  \item gira \(\frac{13\pi}{9}\) radianes en sentido antihorario;
  \item corrige \(95^\circ\) en sentido horario;
  \item se comprueba si la orientacion final cae dentro del sector de seguridad
  \([210^\circ,250^\circ]\).
\end{enumerate}

Determina la orientacion final en grados y en radianes.
\end{exercise}

\begin{shortanswer}{[CTX-C05.S01-01]}
La orientacion final es \(200^\circ\), es decir,
\[
\frac{10\pi}{9}\text{ rad}.
\]
La camara termina fuera del sector de seguridad.
\end{shortanswer}

\begin{fullsolution}{[CTX-C05.S01-01]}
Primero convertimos
\[
\frac{13\pi}{9}\text{ rad}=\frac{13\pi}{9}\cdot \frac{180^\circ}{\pi}=260^\circ.
\]
El giro total es
\[
35^\circ+260^\circ-95^\circ=200^\circ.
\]
Correccion: la suma correcta es
\[
35^\circ+260^\circ=295^\circ,\qquad 295^\circ-95^\circ=200^\circ.
\]
La orientacion final es \(200^\circ\), no \(220^\circ\). En radianes:
\[
200^\circ=\frac{200\pi}{180}=\frac{10\pi}{9}\text{ rad}.
\]
Como \(200^\circ\notin [210^\circ,250^\circ]\), la camara queda fuera del sector de seguridad.
\end{fullsolution}

\section{[C05.S02] Razones trigonometricas en triangulos y poligonos}

\begin{exercise}{[CTX-C05.S02-01] Cubierta de un invernadero}
La seccion frontal de un invernadero es un triangulo isosceles de base \(\SI{14}{m}\). Cada
agua del tejado forma un angulo de \(36^\circ\) con la horizontal. El invernadero tiene una
longitud de \(\SI{9}{m}\).

Calcula:
\begin{enumerate}
  \item la altura de la cumbrera;
  \item la longitud inclinada de cada panel del tejado;
  \item la superficie total de cubierta.
\end{enumerate}
\end{exercise}

\begin{shortanswer}{[CTX-C05.S02-01]}
\[
h=7\tg 36^\circ\approx \SI{5.09}{m},
\qquad
\ell=\frac{7}{\cos 36^\circ}\approx \SI{8.66}{m}.
\]
La superficie total de cubierta es aproximadamente
\[
2\cdot 9\ell \approx \SI{155.8}{m^2}.
\]
\end{shortanswer}

\begin{fullsolution}{[CTX-C05.S02-01]}
La mitad de la base es
\[
\frac{14}{2}=7.
\]
En uno de los triangulos rectangulos formados por la altura:
\[
\tg 36^\circ=\frac{h}{7}
\Longrightarrow
h=7\tg 36^\circ\approx 7\cdot 0.7265\approx \SI{5.09}{m}.
\]
La longitud inclinada del panel es
\[
\cos 36^\circ=\frac{7}{\ell}
\Longrightarrow
\ell=\frac{7}{\cos 36^\circ}\approx \frac{7}{0.8090}\approx \SI{8.66}{m}.
\]
Cada panel tiene area
\[
9\ell \approx 9\cdot 8.66\approx \SI{77.9}{m^2}.
\]
Como hay dos paneles:
\[
A_{total}\approx 2\cdot 77.9=\SI{155.8}{m^2}.
\]
\end{fullsolution}

\section{[C05.S03] Angulos notables y reduccion al primer cuadrante}

\begin{exercise}{[CTX-C05.S03-01] Componentes exactas de un brazo robotico}
Un brazo robotico de longitud \(\SI{12}{m}\) queda orientado con angulo \(210^\circ\) medido
desde el eje \(x\) positivo.

Determina exactamente:
\begin{enumerate}
  \item la componente horizontal;
  \item la componente vertical;
  \item la distancia horizontal del extremo al eje \(y\).
\end{enumerate}
\end{exercise}

\begin{shortanswer}{[CTX-C05.S03-01]}
\[
x=12\cos 210^\circ=-6\sqrt{3},
\qquad
y=12\sen 210^\circ=-6.
\]
La distancia horizontal al eje \(y\) es \(6\sqrt{3}\) metros.
\end{shortanswer}

\begin{fullsolution}{[CTX-C05.S03-01]}
El angulo \(210^\circ\) esta en el tercer cuadrante y su angulo de referencia es \(30^\circ\).
Por tanto,
\[
\cos 210^\circ=-\cos 30^\circ=-\frac{\sqrt{3}}{2},
\qquad
\sen 210^\circ=-\sen 30^\circ=-\frac{1}{2}.
\]
Las componentes del vector posicion son
\[
x=12\cos 210^\circ=12\left(-\frac{\sqrt{3}}{2}\right)=-6\sqrt{3},
\]
\[
y=12\sen 210^\circ=12\left(-\frac{1}{2}\right)=-6.
\]
La distancia horizontal al eje \(y\) es el valor absoluto de la abscisa:
\[
|x|=6\sqrt{3}.
\]
\end{fullsolution}

\section{[C05.S04] Calculo de razones a partir de una razon conocida}

\begin{exercise}{[CTX-C05.S04-01] Sensor angular en el segundo cuadrante}
Un sensor de orientacion registra que el coseno del angulo de un haz vale
\[
\cos\alpha=-\frac{4}{5},
\]
y ademas el haz esta en el segundo cuadrante.

Calcula:
\begin{enumerate}
  \item \(\sen\alpha\) y \(\tg\alpha\);
  \item las coordenadas del extremo del haz si su longitud es \(\SI{25}{m}\).
\end{enumerate}
\end{exercise}

\begin{shortanswer}{[CTX-C05.S04-01]}
\[
\sen\alpha=\frac{3}{5},
\qquad
\tg\alpha=-\frac{3}{4}.
\]
Las coordenadas del extremo son
\[
(-20,15).
\]
\end{shortanswer}

\begin{fullsolution}{[CTX-C05.S04-01]}
Por la identidad fundamental,
\[
\sen^2\alpha=1-\cos^2\alpha=1-\frac{16}{25}=\frac{9}{25}.
\]
Como \(\alpha\) esta en el segundo cuadrante, el seno es positivo:
\[
\sen\alpha=\frac{3}{5}.
\]
Entonces
\[
\tg\alpha=\frac{\sen\alpha}{\cos\alpha}=
\frac{3/5}{-4/5}=-\frac{3}{4}.
\]
Si el haz mide \(25\) m, sus coordenadas son
\[
(25\cos\alpha,25\sen\alpha)=\left(25\cdot -\frac45,25\cdot \frac35\right)=(-20,15).
\]
\end{fullsolution}

\section{[C05.S05] Identidades y transformaciones trigonometricas}

\begin{exercise}{[CTX-C05.S05-01] Dos formulas para el mismo coeficiente}
En un prototipo de compuerta giratoria aparece el coeficiente
\[
K(\alpha)=\frac{\sen 2\alpha}{1+\cos 2\alpha},
\qquad
\alpha\neq \frac{\pi}{2}+k\pi.
\]

Se pide:
\begin{enumerate}
  \item simplificar \(K(\alpha)\);
  \item hallar las posiciones de una vuelta completa en las que \(K(\alpha)=\frac{\sqrt{3}}{3}\).
\end{enumerate}
\end{exercise}

\begin{shortanswer}{[CTX-C05.S05-01]}
\[
K(\alpha)=\tg\alpha.
\]
En \([0,2\pi)\),
\[
\alpha=\frac{\pi}{6}\quad \text{o}\quad \alpha=\frac{7\pi}{6}.
\]
\end{shortanswer}

\begin{fullsolution}{[CTX-C05.S05-01]}
Usamos
\[
\sen 2\alpha=2\sen\alpha\cos\alpha,
\qquad
1+\cos 2\alpha=2\cos^2\alpha.
\]
Entonces
\[
K(\alpha)=\frac{2\sen\alpha\cos\alpha}{2\cos^2\alpha}=\tg\alpha,
\]
siempre que \(\cos\alpha\neq 0\), condicion compatible con la restriccion dada.

Ahora resolvemos
\[
\tg\alpha=\frac{\sqrt{3}}{3}.
\]
El angulo de referencia es \(\frac{\pi}{6}\) y la tangente es positiva en los cuadrantes I y
III. En \([0,2\pi)\) obtenemos
\[
\alpha=\frac{\pi}{6},\qquad \alpha=\pi+\frac{\pi}{6}=\frac{7\pi}{6}.
\]
\end{fullsolution}

\section{[C05.S06] Ecuaciones trigonometricas}

\begin{exercise}{[CTX-C05.S06-01] Fases de mantenimiento en un sistema oscilante}
Un sistema de compensacion entra en modo de mantenimiento cuando el parametro angular
\(\theta\) satisface
\[
2\sen^2\theta-3\sen\theta+1=0.
\]

Determina todas las fases de mantenimiento en el intervalo \([0,2\pi)\).
\end{exercise}

\begin{shortanswer}{[CTX-C05.S06-01]}
\[
\sen\theta=1\quad \text{o}\quad \sen\theta=\frac12.
\]
Por tanto,
\[
\theta=\frac{\pi}{2},\ \frac{\pi}{6},\ \frac{5\pi}{6}.
\]
\end{shortanswer}

\begin{fullsolution}{[CTX-C05.S06-01]}
Hacemos el cambio
\[
u=\sen\theta.
\]
La ecuacion queda
\[
2u^2-3u+1=0.
\]
Factorizando,
\[
(2u-1)(u-1)=0.
\]
Luego
\[
u=\frac12 \quad \text{o}\quad u=1.
\]
Es decir,
\[
\sen\theta=\frac12 \quad \text{o}\quad \sen\theta=1.
\]
En \([0,2\pi)\),
\[
\sen\theta=\frac12 \Longrightarrow \theta=\frac{\pi}{6},\ \frac{5\pi}{6},
\]
\[
\sen\theta=1 \Longrightarrow \theta=\frac{\pi}{2}.
\]
Las fases de mantenimiento son
\[
\theta=\frac{\pi}{6},\ \frac{\pi}{2},\ \frac{5\pi}{6}.
\]
\end{fullsolution}

\section{[C05.S07] Resolucion de triangulos}

\begin{exercise}{[CTX-C05.S07-01] Parcela triangular instrumentada}
Dos linderos de una parcela miden \(\SI{120}{m}\) y \(\SI{95}{m}\), y forman entre si un
angulo de \(58^\circ\).

Calcula:
\begin{enumerate}
  \item la longitud del tercer lado;
  \item el area de la parcela;
  \item el angulo opuesto al lado de \(\SI{95}{m}\).
\end{enumerate}
\end{exercise}

\begin{shortanswer}{[CTX-C05.S07-01]}
El tercer lado mide aproximadamente \(\SI{106.5}{m}\), el area es aproximadamente
\(\SI{4834}{m^2}\) y el angulo opuesto al lado de \(\SI{95}{m}\) es aproximadamente
\(49.2^\circ\).
\end{shortanswer}

\begin{fullsolution}{[CTX-C05.S07-01]}
Por el teorema del coseno,
\[
c^2=120^2+95^2-2\cdot 120\cdot 95\cos 58^\circ.
\]
Sustituyendo,
\[
c^2=14400+9025-22800\cos 58^\circ\approx 23425-12082.15=11342.85.
\]
Por tanto,
\[
c\approx \sqrt{11342.85}\approx \SI{106.5}{m}.
\]

El area es
\[
A=\frac12 ab\sen C=\frac12\cdot 120\cdot 95\sen 58^\circ
\approx 5700\cdot 0.8480\approx \SI{4834}{m^2}.
\]

Para el angulo \(B\) opuesto al lado de \(95\) m, usamos el teorema del seno:
\[
\frac{\sen B}{95}=\frac{\sen 58^\circ}{106.5}.
\]
Entonces
\[
\sen B\approx \frac{95\cdot 0.8480}{106.5}\approx 0.7565.
\]
Como el angulo buscado es el menor de los tres,
\[
B\approx \arcsen(0.7565)\approx 49.2^\circ.
\]
\end{fullsolution}

\section{[C05.S08] Problemas contextualizados de trigonometria}

\begin{exercise}{[CTX-C05.S08-01] Altura de una antena desde dos puntos de la carretera}
Dos puntos de observacion \(A\) y \(B\) estan sobre una misma carretera recta y separados
\(\SI{50}{m}\). Desde \(A\) el angulo de elevacion a la cima de una antena es \(32^\circ\), y
desde \(B\), que esta mas cerca de la base, el angulo es \(47^\circ\).

Calcula:
\begin{enumerate}
  \item la distancia desde \(B\) hasta la base de la antena;
  \item la altura de la antena;
  \item la distancia directa desde \(B\) hasta la cima.
\end{enumerate}
\end{exercise}

\begin{shortanswer}{[CTX-C05.S08-01]}
La distancia desde \(B\) hasta la base es aproximadamente \(\SI{79.9}{m}\), la altura de la
antena es aproximadamente \(\SI{85.7}{m}\) y la visual desde \(B\) hasta la cima mide
aproximadamente \(\SI{117.1}{m}\).
\end{shortanswer}

\begin{fullsolution}{[CTX-C05.S08-01]}
Sea \(x\) la distancia desde \(B\) hasta la base. Entonces desde \(A\) la distancia horizontal
es \(x+50\).

Por trigonometria,
\[
\tg 47^\circ=\frac{h}{x},
\qquad
\tg 32^\circ=\frac{h}{x+50}.
\]
Igualando ambas expresiones de \(h\),
\[
x\tg 47^\circ=(x+50)\tg 32^\circ.
\]
Despejamos:
\[
x(\tg 47^\circ-\tg 32^\circ)=50\tg 32^\circ,
\]
\[
x=\frac{50\tg 32^\circ}{\tg 47^\circ-\tg 32^\circ}
\approx \frac{50\cdot 0.6249}{1.0724-0.6249}\approx \SI{79.9}{m}.
\]
La altura resulta
\[
h=x\tg 47^\circ\approx 79.9\cdot 1.0724\approx \SI{85.7}{m}.
\]
La distancia visual desde \(B\) a la cima es la hipotenusa:
\[
d=\frac{x}{\cos 47^\circ}\approx \frac{79.9}{0.6820}\approx \SI{117.1}{m}.
\]
\end{fullsolution}
